\documentclass[11pt,a4paper]{article}
\usepackage[utf8]{inputenc}
\usepackage[T2A]{fontenc}
\usepackage[russian]{babel}
\usepackage{amsmath,amssymb,booktabs,geometry,hyperref}
\geometry{margin=2.2cm}
\title{Протонный диод В.~Лившица:\\ кинетическая модель, протокол решающего опыта}
\author{Рабочая заметка (по материалам \texttt{Hydrogen\_Thermoel\_Gen.pdf})}
\date{8 ноября 2026}
\begin{document}\maketitle

\section{Устройство}
Одна общая газовая фаза H$_2$, без атомизатора. Слои: (1) поглощающий электрод-катализатор
H$_2\to 2$H$^+ + 2e^-$; (3) протонопроводящая мембрана; (2) выделяющий электрод H$^+ + e^- \to$ H,
покрытый металлическим барьером, пропускающим атомы и задерживающим молекулы. ЭДС по Нернсту
\begin{equation}
E=\frac{RT}{2F}\ln\frac{P_1}{P_2};
\end{equation}
автор предполагает, что барьер поддерживает одностороннюю диффузию при $P_1=P_2$.

\section{Кинетическая модель}
Шесть стадий; каждая подчиняется закону действующих масс с прямой и обратной скоростью,
связанными детальным балансом \emph{своей} стадии. Асимметрия констант разных стадий и сторон произвольна.
\begin{center}
\begin{tabular}{lll}\toprule
стадия & реакция & параметры\\\midrule
1 & H$_2$ + 2$^*\rightleftarrows$ 2H$^*$(1) & $a_1,\ e_1$\\
2 & H$^*$(1) $\rightleftarrows$ H$^+(\psi_1)+e^-(\varphi{=}0)$ & $b_1,\ \varepsilon_p$\\
3 & H$^+(\psi_2)+e^-(\varphi{=}V)\rightleftarrows$ H$^*$(2) & $b_2,\ e_2$\\
4 & H$^*$(2) $\rightleftarrows$ H(барьер) & $k_t,\ e_b$\\
5 & 2H(барьер) $\rightleftarrows$ H$_2$(газ) & $k_d$\\
6 & H + O$_2$(следы) $\to$ H$_2$O & $k_{ox}$\\\bottomrule
\end{tabular}\end{center}
Мембрана --- сопротивление $R_m$, нагрузка $R_L$. Сумма сродств по замкнутому контуру
\begin{equation}
\sum_{\text{стадии}}A_i=\mu_{\mathrm{H_2}}^{\text{вх}}-\mu_{\mathrm{H_2}}^{\text{вых}}=RT\ln\frac{P_{\text{вх}}}{P_{\text{вых}}}-2FV,
\end{equation}
поэтому при $P_{\text{вх}}=P_{\text{вых}}$ и $k_{ox}=0$ стационар с током невозможен при любых константах
(условие Вегшайдера выполняется тождественно). Барьер меняет скорость, но не ЭДС.

\section{Численные результаты}
\begin{enumerate}
\item $P_1=P_2$, без O$_2$, 20 случайных наборов (энергии $\pm3RT$, скорости $10^{-2}\ldots10^{2}$):
$|I|<2\cdot10^{-15}$ --- ток и ЭДС равны нулю.
\item Калибровка: при $P_{\text{вых}}/P_{\text{вх}}=1.01;\,1.1;\,2$ получено $|V|=\tfrac12\ln(P_{\text{вых}}/P_{\text{вх}})$
(в единицах $RT/F$) с точностью $10^{-5}$ --- закон Нернста.
\item Сток O$_2$: $k_{ox}=10^{-4};\,10^{-3};\,10^{-2}$ даёт $V_{oc}\approx0.09;\,0.9;\,8$~мВ при 300~K.
Следы окислителя превращают ячейку в топливный элемент --- главный источник ложного результата.
\end{enumerate}

\section{Протокол решающего опыта}
\begin{itemize}
\item \textbf{Электролит:} Nafion~117 (25--80\,$^\circ$C) или CsH$_2$PO$_4$ ($\sim$240\,$^\circ$C) или
BaZr$_{0.8}$Y$_{0.2}$O$_3$ (400--600\,$^\circ$C, без воды --- предпочтительно).
\item \textbf{Электроды:} поглощающий --- Pt, пористость 30--50\,\%; выделяющий --- Pt, поверх сплошная
плёнка Pd или Pd--Ag 1--5~мкм без пор (гелиевый течеискатель).
\item \textbf{Газ:} одна камера H$_2$ 0.1--1~бар, 6N; O$_2<1$~ppm (геттер, датчик у ячейки).
\item \textbf{Изотермичность:} медный блок, $\Delta T<0.01$~K между электродами.
\item \textbf{Измерения:} ЭДС электрометром ($\sim$1~мкВ, вход $\ge10^{12}$~Ом), ток КЗ, шум и импеданс:
$T_{\text{эфф}}(f)=S_V/(4k\,\mathrm{Re}Z)$.
\item \textbf{Контроли:} калибровка Нернстом ($P_2/P_1=1.01\Rightarrow0.13$~мВ при 300~K);
переворот ячейки (смена знака); симметричная ячейка (эффект исчезает); дозированный O$_2$ 1--100~ppm;
баланс заряда $\int I\,dt/2F$ против расхода H$_2$ и наработки воды.
\end{itemize}
Мощность, \emph{если бы} ЭДС $E$ существовала, при $R\approx0.1$~Ом$\cdot$см$^2$:
$E^2/4R\approx0.025;\ 2.5;\ 250$~Вт/м$^2$ при $E=1;\ 10;\ 100$~мВ.

\section{Статья 2024 г.}
«Известия АЭН РФ», вып.~1(28), 2024, с.~59--80 («транзитный осмос»). Полный текст в открытом доступе
не найден (на \url{http://www.acelsc.ru/node/21} есть только список выпусков); разбор будет дополнен
после получения PDF.

\section{Вывод}
В модели с произвольной асимметрией стадий ЭДС при $P_1=P_2$ строго ноль; разность давлений даёт точно
Нернста; следы O$_2$ --- милливольты. Опыт простой и чувствительный: ненулевая устойчивая ЭДС при
O$_2<1$~ppm, $\Delta T<0.01$~K, со сменой знака при перевороте и отсутствием в симметричной ячейке была бы
настоящим результатом.
\end{document}
